\subsection{prime-register 可逆化机制的刚性、位复杂度与合成闭包}\label{subsec:conclusion-prime-register-rigidity-complexity-conjugacy-composition}
沿用定义 \ref{def:pom-prime-fiber} 与定义 \ref{def:pom-fold-prime-lift} 的记号，记
\[
\mathcal P=\NN^{(\mathbb P)},\qquad
N(r)=\prod_{p\in\mathbb P}p^{r(p)},
\]
并以命题 \ref{prop:pom-fold-prime-lift-injective} 作为可逆化基底。
下述命题组把时间寻址、整数化代价、策略变化与投影合成统一写成同一条可审计代数链。

\begin{theorem}[时间寻址 Gödel 证书的刚性与唯一性]\label{thm:conclusion-prime-time-addressed-godel-rigidity}
令
\[
\mathcal A:=\NN^{(\NN)}=\left\{a=(a_t)_{t\ge 1}\mid a_t=0\ \text{仅除有限多个 }t\right\},
\]
并定义移位算子 \(\mathrm{sh}:\mathcal A\to\mathcal A\) 为
\[
(\mathrm{sh}(a))_1=0,\qquad (\mathrm{sh}(a))_{t+1}=a_t.
\]
设 \(\Gamma:\mathcal A\to\NN_{\ge 1}\) 满足存在两两不同素数列 \((q_t)_{t\ge 1}\) 使
\[
\Gamma(a)=\prod_{t\ge 1}q_t^{a_t}\qquad(a\in\mathcal A),
\]
并且存在乘法幺半群自同态 \(\Sigma:\NN_{\ge 1}\to\NN_{\ge 1}\) 使
\[
\Gamma(\mathrm{sh}(a))=\Sigma(\Gamma(a)),\qquad \forall a\in\mathcal A.
\]
则成立：
\begin{enumerate}
  \item 在由 \(\{q_t\}_{t\ge 1}\) 生成的子幺半群上，\(\Sigma\) 必满足
  \[
  \Sigma(q_t)=q_{t+1},\qquad
  \Sigma\!\left(\prod_{t\ge 1}q_t^{\alpha_t}\right)=\prod_{t\ge 1}q_{t+1}^{\alpha_t}.
  \]
  \item 若进一步假设 \(\Sigma\) 是 \(\NN_{\ge 1}\) 的自同构，则存在唯一素数置换
  \(\pi:\mathbb P\to\mathbb P\) 使
  \[
  \Sigma\!\left(\prod_{p\in\mathbb P}p^{\alpha(p)}\right)
  =
  \prod_{p\in\mathbb P}\pi(p)^{\alpha(p)},
  \]
  且在时间寻址轴上 \(\pi(q_t)=q_{t+1}\)。
\end{enumerate}
\end{theorem}
\begin{proof}
取 \(e_t\in\mathcal A\) 为第 \(t\) 位为 \(1\) 其余为 \(0\) 的序列。
由移位相容性
\[
\Sigma(q_t)
=\Sigma(\Gamma(e_t))
=\Gamma(\mathrm{sh}(e_t))
=\Gamma(e_{t+1})
=q_{t+1},
\]
故第一式成立。
对任意有限支撑指数列 \((\alpha_t)\)，由 \(\Sigma\) 乘法同态性质得
\[
\Sigma\!\left(\prod_{t\ge 1}q_t^{\alpha_t}\right)
=\prod_{t\ge 1}\Sigma(q_t)^{\alpha_t}
=\prod_{t\ge 1}q_{t+1}^{\alpha_t},
\]
即得第 (1) 条。

若 \(\Sigma\) 为自同构，则由算术基本定理，\(\Sigma\) 在素数上的取值必须构成素数集的置换 \(\pi\)，并唯一决定 \(\Sigma\) 在整个 \(\NN_{\ge 1}\) 上的作用：
\[
\Sigma\!\left(\prod_{p}p^{\alpha(p)}\right)=\prod_p\pi(p)^{\alpha(p)}.
\]
结合上式 \(\Sigma(q_t)=q_{t+1}\) 即得 \(\pi(q_t)=q_{t+1}\)。
唯一性由唯一分解直接推出。证毕。
\end{proof}

\begin{theorem}[prime-register 整数化的超线性位复杂度下界]\label{thm:conclusion-prime-integerization-superlinear-bitlength}
固定确定性策略 \(\mathsf S\)，并沿用定义 \ref{def:pom-fold-prime-lift} 中的事件序列
\((e_t(\omega))_{t=1}^{T_{m+1}^{\mathsf S}(\omega)}\) 与编码 \(\mathrm{code}(e_t)\in\NN_{\ge 1}\)。
定义 Gödel 整数
\[
G_{\mathsf S}(\omega):=N\!\bigl(r_{\mathsf S}(\omega)\bigr)
=\prod_{t=1}^{T_{m+1}^{\mathsf S}(\omega)}p_t^{\,\mathrm{code}(e_t(\omega))}.
\]
则对任意 \(\omega\) 有
\[
\log_2 G_{\mathsf S}(\omega)
\ge
\log_2\!\bigl((T_{m+1}^{\mathsf S}(\omega)+1)!\bigr)
\ge
\frac{(T_{m+1}^{\mathsf S}(\omega)+1)\log(T_{m+1}^{\mathsf S}(\omega)+1)-(T_{m+1}^{\mathsf S}(\omega)+1)+1}{\ln 2}.
\]
特别地，
\[
\log_2 G_{\mathsf S}(\omega)=\Omega\!\bigl(T_{m+1}^{\mathsf S}(\omega)\log T_{m+1}^{\mathsf S}(\omega)\bigr).
\]
在均匀基线 \(\omega\sim\mathrm{Unif}(\Omega_{m+1})\) 下，存在常数 \(C>0\) 使
\[
\EE[\log_2 G_{\mathsf S}(\omega)]
\ge
\frac{4}{27\ln 2}\,m\log m-Cm.
\]
\end{theorem}
\begin{proof}
由定义 \ref{def:pom-fold-prime-lift}，\(\mathrm{code}(e_t)\ge 1\)，故
\[
G_{\mathsf S}(\omega)\ge \prod_{t=1}^{T}p_t,\qquad T:=T_{m+1}^{\mathsf S}(\omega).
\]
又 \(p_t\ge t+1\)，于是
\[
G_{\mathsf S}(\omega)\ge \prod_{t=1}^{T}(t+1)=(T+1)!.
\]
取对数得第一不等式。
再由
\[
\log (T+1)!
=\sum_{k=1}^{T+1}\log k
\ge
\int_1^{T+1}\log x\,\mathrm dx
=(T+1)\log(T+1)-(T+1)+1
\]
得到第二不等式，从而点态 \(\Omega(T\log T)\) 成立。

令
\[
h(t):=(t+1)\log(t+1)-(t+1),\qquad t\ge 0.
\]
有 \(h''(t)=1/(t+1)>0\)，故 \(h\) 凸。
由上界链与 Jensen 不等式，
\[
\EE[\log_2 G_{\mathsf S}]
\ge
\frac{1}{\ln 2}\,\EE[h(T)]
\ge
\frac{1}{\ln 2}\,h(\EE[T]).
\]
再由推论 \ref{cor:pom-prime-residual-linear-lb}，
\[
\liminf_{m\to\infty}\frac{\EE[T]}{m}\ge \frac{4}{27},
\]
代回即得
\[
\EE[\log_2 G_{\mathsf S}]
\ge
\frac{4}{27\ln 2}\,m\log m-Cm
\]
（某常数 \(C>0\)）。证毕。
\end{proof}

\begin{theorem}[全息—深度分离原理：低前缀信息与任意可计算时间深度可共存]\label{thm:conclusion-holography-logical-depth-separation}
固定一台前缀自由通用机 \(U\)，并以 \(K(\cdot)\) 表示前缀 Kolmogorov 复杂度。
令 \(g:\NN\to\NN\) 为任意全可计算、单调不减且无界的函数。
则存在一个可计算无穷二进制序列 \(X\in\{0,1\}^{\NN}\)，以及一列前缀长度
\[
N_m:=\frac{m(m+1)}{2}\qquad(m\ge 1),
\]
使得对所有充分大的 \(m\) 同时成立：
\begin{enumerate}
  \item 存在常数 \(C_1>0\) 使
  \[
  K\!\bigl(X\upharpoonright N_m\bigr)\ \le\ C_1\log m;
  \]
  \item 令 \(K^t(\sigma)\) 为时间界 \(t\) 下的时间有界前缀复杂度，并定义 Bennett 型逻辑深度
  \[
  \mathrm{depth}_s(\sigma):=\min\{t:\ K^t(\sigma)\le K(\sigma)+s\}.
  \]
  则存在常数 \(C_2>0\) 使
  \[
  \mathrm{depth}_{C_2\log m}\!\bigl(X\upharpoonright N_m\bigr)\ >\ g(m).
  \]
\end{enumerate}
\end{theorem}
\begin{proof}
取常数 \(c\) 并令
\[
L(m):=\bigl\lceil 2\log_2 m\bigr\rceil+c.
\]
对每个 \(m\) 定义集合
\[
\mathcal S_m:=
\Bigl\{
\tau\in\{0,1\}^m:\ \exists\,p,\ |p|\le L(m)\ \text{且}\ U(p)\ \text{在}\ \le g(m)\ \text{步内输出}\ \tau
\Bigr\}.
\]
由于 \(U\) 前缀自由，长度 \(\le L(m)\) 的程序数至多为 \(2^{L(m)}\)，故
\(
|\mathcal S_m|\le 2^{L(m)}\ll m^2
\)。
当 \(m\) 充分大时有 \(|\mathcal S_m|<2^m\)，从而 \(\{0,1\}^m\setminus \mathcal S_m\neq\varnothing\)。
令
\[
B_m:=\min\nolimits_{\mathrm{lex}}\bigl(\{0,1\}^m\setminus \mathcal S_m\bigr)
\]
为字典序最小的缺失块。
由于 \(g(m)\) 为可计算且时间界有限，可通过有限枚举全部 \(|p|\le L(m)\) 的程序并在步数 \(\le g(m)\) 内模拟判定 \(U(p)\) 的输出，从而 \(\mathcal S_m\) 与 \(B_m\) 均可计算。

定义
\[
X:=B_1B_2B_3\cdots .
\]
于是 \(X\) 可计算。
并且由给定 \(m\) 可逐块计算 \(B_1,\dots,B_m\) 并输出其拼接 \(X\upharpoonright N_m\)，故存在常数 \(C_1\) 使
\(
K(X\upharpoonright N_m)\le C_1\log m
\)。

设反之存在程序 \(p\) 满足
\[
|p|\le K(X\upharpoonright N_m)+C_2\log m,
\qquad
U(p)\ \text{在}\ \le g(m)\ \text{步内输出}\ X\upharpoonright N_m.
\]
构造程序 \(q\)：将 \(m\) 与 \(p\) 作为内嵌常量，模拟 \(U(p)\) 的输出并丢弃前 \(N_{m-1}\) 位，仅输出接下来的 \(m\) 位后停机。
则 \(U(q)\) 在 \(\le g(m)\) 步内输出 \(B_m\)。
并且 \(|q|\le |p|+O(\log m)\)。
取 \(C_2\) 足够大，使得对充分大的 \(m\) 有 \(|q|\le L(m)\)，则 \(B_m\in\mathcal S_m\)，与 \(B_m\notin\mathcal S_m\) 矛盾。
故上述 \(p\) 不存在，从而
\(
\mathrm{depth}_{C_2\log m}(X\upharpoonright N_m)>g(m)
\)
对充分大的 \(m\) 成立。证毕。
\end{proof}

\begin{corollary}[信息饱和而时间可任意爆炸：以 \(N_m\) 为尺度的表述]\label{cor:conclusion-holography-logical-depth-separation-Nm}
存在可计算序列 \(X\) 与无穷前缀长度列 \((N_m)\) 使
\[
K\!\bigl(X\upharpoonright N_m\bigr)=O(\log N_m),
\]
但对任意给定可计算无界单调函数 \(g\)，仍可使
\[
\mathrm{depth}_{O(\log N_m)}\!\bigl(X\upharpoonright N_m\bigr)\ \ge\ g(\sqrt{N_m})
\]
最终成立。
\end{corollary}
\begin{proof}
由定理 \ref{thm:conclusion-holography-logical-depth-separation} 取 \(N_m=\frac{m(m+1)}2\asymp m^2\)，并以 \(\log m\asymp \log N_m\) 代入即可。证毕。
\end{proof}

\begin{remark}[低信息并不蕴含可快速识别]\label{rem:conclusion-low-information-not-fast-recognition}
对定理 \ref{thm:conclusion-holography-logical-depth-separation} 的 \(X\) 与充分大的 \(m\)，不等式
\(
\mathrm{depth}_{C_2\log m}(X\upharpoonright N_m)>g(m)
\)
等价于断言：任意长度不超过 \(K(X\upharpoonright N_m)+C_2\log m\) 的程序若输出 \(X\upharpoonright N_m\)，其运行时间必超过 \(g(m)\)。
因此，概念描述的对数级短性与相干输出的不可规避计算成本可在同一可计算对象中严格共存。
\end{remark}

\begin{theorem}[策略变化下 prime-register 证书系统的可计算共轭]\label{thm:conclusion-prime-strategy-computable-conjugacy}
设 \(\mathsf S,\mathsf S'\) 为两种确定性策略，并定义
\[
\widetilde{\Fold}_m^{\mathsf S},\ \widetilde{\Fold}_m^{\mathsf S'}:\Omega_m\to X_m\times\mathcal P
\]
如定义 \ref{def:pom-fold-prime-lift}。
令
\[
\mathcal I_m^{\mathsf S}:=\widetilde{\Fold}_m^{\mathsf S}(\Omega_m),\qquad
\mathcal I_m^{\mathsf S'}:=\widetilde{\Fold}_m^{\mathsf S'}(\Omega_m),
\]
并以命题 \ref{prop:pom-fold-prime-lift-injective} 的逆过程记号定义
\[
\Psi_{\mathsf S\to\mathsf S'}
:=
\widetilde{\Fold}_m^{\mathsf S'}\circ
\mathrm{UNLIFT}_m^{\mathsf S}
:\mathcal I_m^{\mathsf S}\to\mathcal I_m^{\mathsf S'}.
\]
则：
\begin{enumerate}
  \item \(\Psi_{\mathsf S\to\mathsf S'}\) 为双射，且逆为 \(\Psi_{\mathsf S'\to\mathsf S}\)；
  \item 若 \(\Psi_{\mathsf S\to\mathsf S'}(x,r)=(x',r')\)，则 \(x'=x\)；
  \item 在 \(\Omega_m\) 上成立共轭关系
  \[
  \widetilde{\Fold}_m^{\mathsf S'}
  =
  \Psi_{\mathsf S\to\mathsf S'}\circ \widetilde{\Fold}_m^{\mathsf S}.
  \]
\end{enumerate}
\end{theorem}
\begin{proof}
任取 \((x,r)\in \mathcal I_m^{\mathsf S}\)，令
\[
\omega:=\mathrm{UNLIFT}_m^{\mathsf S}(x,r).
\]
则
\[
\Psi_{\mathsf S\to\mathsf S'}(x,r)=\widetilde{\Fold}_m^{\mathsf S'}(\omega)\in \mathcal I_m^{\mathsf S'},
\]
故映射良定义。
同理
\[
\Psi_{\mathsf S'\to\mathsf S}\circ\Psi_{\mathsf S\to\mathsf S'}(x,r)
=
\widetilde{\Fold}_m^{\mathsf S}\!\left(
\mathrm{UNLIFT}_m^{\mathsf S'}\!\left(\widetilde{\Fold}_m^{\mathsf S'}(\omega)\right)
\right)
=
\widetilde{\Fold}_m^{\mathsf S}(\omega)
=(x,r),
\]
另一方向同理，故互逆，得第 (1) 条。

由定义
\[
\widetilde{\Fold}_m^{\mathsf S'}(\omega)=\bigl(\Fold_m(\omega),\,r_{\mathsf S'}(\omega)\bigr),
\]
而 \((x,r)\in\mathcal I_m^{\mathsf S}\) 意味着 \(x=\Fold_m(\omega)\)，故 \(x'=x\)，得第 (2) 条。

对任意 \(\omega\in\Omega_m\)，
\[
\Psi_{\mathsf S\to\mathsf S'}\!\bigl(\widetilde{\Fold}_m^{\mathsf S}(\omega)\bigr)
=
\widetilde{\Fold}_m^{\mathsf S'}\!\left(
\mathrm{UNLIFT}_m^{\mathsf S}\!\left(\widetilde{\Fold}_m^{\mathsf S}(\omega)\right)
\right)
=
\widetilde{\Fold}_m^{\mathsf S'}(\omega),
\]
即得第 (3) 条。证毕。
\end{proof}

\begin{theorem}[投影合成下的素数交错分配律]\label{thm:conclusion-prime-interleaving-composition-law}
设 \(f:\Omega\to X\)、\(g:X\to Y\) 为两级投影。
假设存在确定性策略 \(\mathsf S_f,\mathsf S_g\) 及寄存器读出
\[
\widetilde f^{\mathsf S_f}(\omega):=\bigl(f(\omega),r_f(\omega)\bigr)\in X\times\mathcal P,\qquad
\widetilde g^{\mathsf S_g}(x):=\bigl(g(x),r_g(x)\bigr)\in Y\times\mathcal P,
\]
并且二者均在对象层单射且可由寄存器回放逆过程重建。
定义交错嵌入 \(I_{\mathrm{odd}},I_{\mathrm{even}}:\mathcal P\to\mathcal P\)：
\[
I_{\mathrm{odd}}(r)(p_{2t-1})=r(p_t),\quad I_{\mathrm{odd}}(r)(p_{2t})=0,
\]
\[
I_{\mathrm{even}}(r)(p_{2t})=r(p_t),\quad I_{\mathrm{even}}(r)(p_{2t-1})=0.
\]
令合成寄存器
\[
r_{g\circ f}(\omega)
:=
I_{\mathrm{even}}(r_f(\omega))
+
I_{\mathrm{odd}}(r_g(f(\omega))),
\]
并定义
\[
\widetilde{(g\circ f)}^{\mathrm{mix}}(\omega)
:=
\bigl(g(f(\omega)),\,r_{g\circ f}(\omega)\bigr).
\]
则 \(\widetilde{(g\circ f)}^{\mathrm{mix}}\) 为单射，且可由 \((y,r_{g\circ f})\) 确定性重建 \(\omega\)。
\end{theorem}
\begin{proof}
由定义，\(I_{\mathrm{odd}}(\mathcal P)\) 与 \(I_{\mathrm{even}}(\mathcal P)\) 支撑不交，故可从 \(r_{g\circ f}\) 唯一分离
\[
r_f^\sharp(p_t):=r_{g\circ f}(p_{2t}),\qquad
r_g^\sharp(p_t):=r_{g\circ f}(p_{2t-1})\qquad (t\ge 1).
\]
先用 \((y,r_g^\sharp)\) 经 \(\widetilde g^{\mathsf S_g}\) 的逆过程恢复唯一 \(x\in X\)；
再用 \((x,r_f^\sharp)\) 经 \(\widetilde f^{\mathsf S_f}\) 的逆过程恢复唯一 \(\omega\in\Omega\)。
故 \((y,r_{g\circ f})\) 唯一决定 \(\omega\)，即 \(\widetilde{(g\circ f)}^{\mathrm{mix}}\) 单射且可回放。证毕。
\end{proof}
